% Generated by report_gen.py (ai-ee v1) - do not hand-edit.
\documentclass[11pt,a4paper]{article}
\usepackage[margin=2.2cm]{geometry}
\usepackage{graphicx}
\usepackage{booktabs}
\usepackage{longtable}
\usepackage{pdfpages}
\makeatletter
\@ifundefined{KV@Gin@artifact}{\define@key{Gin}{artifact}[]{}}{}
\makeatother
\usepackage{xcolor}
\usepackage[hidelinks]{hyperref}
\setcounter{tocdepth}{1}
\begin{document}
\sloppy
\begin{center}
{\LARGE\bfseries rf-de-20m --- Design Document}\\[6pt]
{\large ai-ee v1 pipeline}\\[2pt]
generated 2026-08-07 22:06:16
\end{center}
\tableofcontents

\section{Board and Run Metadata}
{\small
\begin{longtable}{lp{11cm}}
\toprule
\textbf{Field} & \textbf{Value} \\
\midrule
\endhead
board & rf-de-20m \\
workspace & \texttt{boards/rf-de-20m} \\
phase & P2 \\
gate status & no gates recorded yet \\
state created & 2026-08-07T20:58:16 \\
state updated & 2026-08-07T21:33:52 \\
\bottomrule
\end{longtable}
}
\section{Overview}
\subsection*{BRIEF - rf-de-20m: 20 MHz Class E GaN RF power stage, 200 W}
Owner-approved 2026-08-07. Full derivation + rejected alternatives: \texttt{LEARNINGS.md} (this dir). \textbf{Read that file before changing any number here} - the operating point is frozen and the Class DE / gate-drive-transformer alternatives are already closed.

\subsubsection*{Scope (deliberately narrow - owner asked for power stage ONLY)}
IN scope: gate driver, GaN power stage, output filter + impedance match, 5 V housekeeping. OUT of scope: no MCU, no PWM generation, no telemetry, no protection logic, no enclosure. PWM arrives from an external signal generator via SMA. This is a bench power stage.

\subsubsection*{Hard constraints}
\begin{itemize}
\item \textbf{100\% JLCPCB PCBA - zero off-catalog parts.} Owner ruled on this explicitly; it is the reason
the topology is Class E and not the originally requested Class DE (no LCSC-stocked driver can
switch a 20 MHz half-bridge; see LEARNINGS). Every part must be LCSC/JLC stocked.
\item 4-layer PCB. L2 = unbroken GND directly under the power loop. Non-negotiable at 20 MHz.
\item Minimise spend by reducing SCOPE, never by reducing quality.
\end{itemize}
\subsubsection*{Operating point (frozen - do not re-derive)}
\begin{quote}\small\ttfamily\raggedright
\textbar{} Param \textbar{} Value \textbar{}\\
\textbar{}---\textbar{}---\textbar{}\\
\textbar{} Frequency \textbar{} 20 MHz \textbar{}\\
\textbar{} Output power \textbar{} 200 W into 50 ohm \textbar{}\\
\textbar{} Bus Vdd \textbar{} 40 V DC (needs \textgreater{}=6 A supply) \textbar{}\\
\textbar{} R\_opt (Class E) \textbar{} 4.614 ohm \textbar{}\\
\textbar{} Vds peak \textbar{} 142.5 V (3.562 x Vdd), 1.4x margin on 200 V part \textbar{}\\
\textbar{} Ids peak \textbar{} \textasciitilde{}16 A \textbar{}\\
\textbar{} I\_dc \textbar{} \textasciitilde{}5.8 A \textbar{}\\
\textbar{} Expected efficiency \textbar{} {*}{*}80-85\%{*}{*} (NOT 86\% - see thermal note) \textbar{}
\end{quote}
\subsubsection*{Core parts (both LCSC-verified in stock)}
\begin{itemize}
\item \textbf{Switch:} EPC2019 eGaN FET - 200 V, 36 mohm, Coss 110 pF, Qoss 18 nC@100 V. \textbf{LCSC C2836675},
\textasciitilde{}\$2.17. 1.35 mm chip-scale LGA, bottom-cooled through its solder balls.
\item \textbf{Driver:} LMG1020YFFR - 5 V, 7 A source / 5 A sink, 3 ns prop delay, 1 ns min pulse, 60 MHz.
\textbf{LCSC C6423790}, \textasciitilde{}\$0.36, 0.8x1.2 mm WCSP. NB: LCSC brands it "Tokmas", not TI - flag for
authenticity check on receipt.
\item \textbf{5 V rail:} small 40 V-input buck. NOT an LDO - gate charge draw is Qg{*}f \textasciitilde{} 0.1 A, so an LDO
would burn (40-5){*}0.1 = 3.5 W. Pick an LCSC-stocked 40 V-capable buck.
\end{itemize}
\subsubsection*{Output network (Class E, Q\_L = 5)}
Drain --+-- C\_shunt --GND        C\_shunt 317 pF TOTAL (Coss 110 pF + \textasciitilde{}200 pF ext C0G)

\begin{quote}\small\ttfamily\raggedright
            \textbar{}
\end{quote}
+-- L\_s -- C\_s --+-- L\_m --+-- SMA out (50 ohm)

\begin{quote}\small\ttfamily\raggedright
                             \textbar{}         \textbar{}
\end{quote}
(--)      C\_m -- GND

L\_s 184 nH \textbar{} C\_s \textasciitilde{}447 pF \textbar{} L\_m \textasciitilde{}115 nH \textbar{} C\_m \textasciitilde{}500 pF (match Q 3.14, 4.614 -\textgreater{} 50 ohm)

\subsubsection*{MUST-RESOLVE during build (do not silently skip)}
1. \textbf{Tank current is 6.6 A rms.} In a series-resonant tank the circulating current equals the load current - loaded Q sets VOLTAGE across L/C, not current. Size L and C by \emph{dissipation}, not current rating. 2. \textbf{The output inductor is the hottest component on the board} (\textasciitilde{}7 W at Q=150), hotter than the GaN. If no LCSC RF inductor carries 6.6 A rms at high Q: parallel several (splits current AND ESR), or use a PCB air-core spiral (free, JLC-native, costs area). 3. \textbf{Parallel several C0G caps} for C\_s / C\_shunt to split ripple current. L/C see \textasciitilde{}215 V peak -\textgreater{} specify \textgreater{}=250 V, ideally 500 V C0G. No X7R anywhere in the tank (voltage/temp coefficient). 4. \textbf{Verify LMG1020 input} accepts a 50 ohm-terminated logic-level SMA signal directly; add bias / a fast buffer if not. 5. Total dissipation 35-50 W. Thermal path for the EPC2019 is vias straight into plane copper.

\subsubsection*{Layout rules}
\begin{itemize}
\item Gate loop (LMG1020 -\textgreater{} EPC2019 gate -\textgreater{} source) must be the tightest loop on the board. The WCSP
driver exists to be placed millimetres from the FET; do not waste that.
\item Power loop area, not trace width, sets the ringing that the 1.4x Vds derate protects.
\item Skin depth in Cu at 20 MHz is 14.6 um; 1 oz = 35 um is only \textasciitilde{}2.4 skin depths. Pour tank and
power-loop conductors wide - do NOT size by DC current.
\item Decoupling: multiple small-case C0G right at the drain/source. Bulk electrolytic is useless at
20 MHz - it is there only for the 5.8 A DC feed.
\item RF choke feeding the drain must be a real RF choke, self-resonance well above 20 MHz.
\item Keep the 50 ohm output trace controlled-impedance to the SMA.
\end{itemize}
\section{Requirements}
\subsection*{requirements.md - rf-de-20m}
20 MHz Class E GaN RF power stage, 200 W into 50 ohm. Source: \texttt{brief/BRIEF.md} + \texttt{brief/RESEARCH-LEARNINGS.md}, owner-approved 2026-08-07.

\textbf{Frozen inputs - not open for re-derivation.} The topology, operating point and the two core parts below were decided by the owner after a full costed trade study. They are recorded here as settled requirements, not as proposals. Class DE was evaluated and REJECTED (no LCSC-stocked gate driver can switch a 20 MHz half-bridge). Do not reopen.

---

\subsubsection*{1. Function}
A bench RF power amplifier stage: a single-ended Class E inverter that converts a 40 V DC bus into 200 W of 20 MHz RF power delivered to a 50 ohm load. An externally generated 20 MHz PWM drive signal arrives on an SMA input, is buffered by a low-side GaN gate driver, and switches a single ground-referenced eGaN FET. The FET's drain feeds a shunt capacitance plus a series resonant tank and an L-match that transforms the 4.614 ohm Class E load line up to 50 ohm at the output SMA. Scope is the POWER STAGE ONLY. There is no MCU, no on-board PWM generation, no telemetry, and no protection logic on this board.

\subsubsection*{2. Interfaces}
\begin{quote}\small\ttfamily\raggedright
\textbar{} \# \textbar{} Interface \textbar{} Direction \textbar{} Standard / detail \textbar{} Connector \textbar{}\\
\textbar{}---\textbar{}---\textbar{}---\textbar{}---\textbar{}---\textbar{}\\
\textbar{} I1 \textbar{} RF drive input \textbar{} in \textbar{} 20 MHz logic-level PWM from an external signal generator; 50 ohm terminated at the board. Feeds the LMG1020 input. Adjustable duty cycle expected (\textasciitilde{}50\% nominal for Class E). \textbar{} SMA (stated in brief) \textbar{}\\
\textbar{} I2 \textbar{} RF power output \textbar{} out \textbar{} 20 MHz, 200 W CW into 50 ohm. Controlled-impedance 50 ohm trace from the L-match to the connector. \textbar{} SMA (stated in brief) - see Q10 \textbar{}\\
\textbar{} I3 \textbar{} DC bus input \textbar{} in \textbar{} 40 V DC, \textgreater{}=6 A. \textbar{} NOT STATED - see Q5 \textbar{}\\
\textbar{} I4 \textbar{} 5 V housekeeping \textbar{} in or internal \textbar{} 5 V rail for the LMG1020. Source not yet decided: onboard 40 V-input buck vs external 5 V feed. \textbar{} see Q9 \textbar{}
\end{quote}
Notes carried from the brief that bind the interfaces:

\begin{itemize}
\item The 50 ohm output trace must be controlled-impedance all the way to the output connector.
\item MUST-RESOLVE at design time: verify the LMG1020 input accepts a 50 ohm-terminated logic-level
SMA signal directly; add bias or a fast buffer if it does not.
\item No other external interfaces exist. No status LEDs, no buttons, no debug/telemetry header are
required by the brief. ASSUMED: none wanted (scope is deliberately bare bones).
\end{itemize}
\subsubsection*{3. Power}
\textbf{Main bus (frozen):}

\begin{quote}\small\ttfamily\raggedright
\textbar{} Param \textbar{} Value \textbar{} Status \textbar{}\\
\textbar{}---\textbar{}---\textbar{}---\textbar{}\\
\textbar{} Bus Vdd \textbar{} 40 V DC \textbar{} frozen \textbar{}\\
\textbar{} Bus current, DC average \textbar{} \textasciitilde{}5.8 A at 85\% efficiency \textbar{} frozen \textbar{}\\
\textbar{} Supply requirement \textbar{} \textgreater{}=6 A capable at 40 V \textbar{} frozen \textbar{}\\
\textbar{} Output power \textbar{} 200 W into 50 ohm \textbar{} frozen \textbar{}\\
\textbar{} Class E load line R\_opt \textbar{} 4.614 ohm \textbar{} frozen \textbar{}\\
\textbar{} Vds peak \textbar{} 142.5 V (3.562 x Vdd) -\textgreater{} 1.4x margin on the 200 V EPC2019 \textbar{} frozen \textbar{}\\
\textbar{} Ids peak \textbar{} \textasciitilde{}16 A \textbar{} frozen \textbar{}\\
\textbar{} Tank current \textbar{} 6.6 A rms (equals load current - loaded Q sets VOLTAGE across L/C, not current) \textbar{} frozen \textbar{}\\
\textbar{} Expected efficiency \textbar{} 80-85\% (NOT 86\%) \textbar{} frozen \textbar{}\\
\textbar{} Total dissipation \textbar{} 35-50 W, most of it in the passives, not the semiconductor \textbar{} frozen \textbar{}
\end{quote}
\textbf{Housekeeping rail:}

\begin{itemize}
\item 5 V for the LMG1020 gate driver. Must be a buck, NOT an LDO: gate-charge draw is Qg{*}f \textasciitilde{} 0.1 A,
so a 40 V -\textgreater{} 5 V LDO would burn (40-5){*}0.1 = 3.5 W.
\item GUESS: 5 V rail budget 150-250 mA (0.1 A gate charge + driver quiescent + margin). To be
confirmed against the actual EPC2019 Qg at the chosen drive voltage during design.
\item If an onboard buck is used it must be LCSC-stocked and rated for a 40 V input (see Q9).
\end{itemize}
\textbf{Frozen part selections:}

\begin{itemize}
\item Switch: EPC2019 eGaN FET, 200 V / 36 mohm / Coss 110 pF / Qoss 18 nC @ 100 V. LCSC C2836675,
\textasciitilde{}\$2.17. 1.35 mm chip-scale LGA, bottom-cooled through its solder balls.
\item Driver: LMG1020YFFR, 5 V, 7 A source / 5 A sink, 3 ns prop delay, 1 ns min pulse, 60 MHz.
LCSC C6423790, \textasciitilde{}\$0.36, 0.8 x 1.2 mm WCSP. NB: LCSC brands it "Tokmas", not TI - flag for an
authenticity check on receipt.
\end{itemize}
\textbf{Frozen output network (Class E, Q\_L = 5):}

Drain --+-- C\_shunt --GND        C\_shunt 317 pF TOTAL (Coss 110 pF + \textasciitilde{}200 pF external C0G)

\begin{quote}\small\ttfamily\raggedright
            \textbar{}
\end{quote}
+-- L\_s -- C\_s --+-- L\_m --+-- SMA out (50 ohm)

\begin{quote}\small\ttfamily\raggedright
                                       \textbar{}
\end{quote}
C\_m -- GND

L\_s 184 nH \textbar{} C\_s \textasciitilde{}447 pF \textbar{} L\_m \textasciitilde{}115 nH \textbar{} C\_m \textasciitilde{}500 pF (match Q 3.14, 4.614 -\textgreater{} 50 ohm)

Binding design constraints on that network (from the brief's MUST-RESOLVE list):

\begin{itemize}
\item The output inductor is the hottest component on the board (\textasciitilde{}7 W at Q=150), hotter than the GaN.
Size L and C by DISSIPATION and Q, not by current rating. If no LCSC RF inductor carries 6.6 A
rms at high Q: parallel several (splits current AND ESR), or use a PCB air-core spiral
(JLC-native, free in BOM, costs board area).
\item Parallel several C0G caps for C\_s and C\_shunt to split ripple current. Tank L/C see \textasciitilde{}215 V peak
-\textgreater{} specify \textgreater{}=250 V, ideally 500 V C0G. No X7R anywhere in the tank.
\item The drain-feed RF choke must be a real RF choke with self-resonance well above 20 MHz.
\item Decoupling: multiple small-case (0402/0603) C0G right at drain/source. Bulk electrolytic is
useless at 20 MHz and exists only to support the 5.8 A DC feed.
\end{itemize}
\subsubsection*{4. Environment}
NOT STATED in the brief. What is known:

\begin{itemize}
\item Bench instrument, no enclosure (explicitly out of scope). Ingress protection: none required.
\item Vibration: none - stationary bench use. ASSUMED.
\item Ambient temperature: not stated - see Q7.
\item Cooling: 35-50 W of dissipation must leave the board. The EPC2019 is bottom-cooled through its
solder balls, so its thermal path is vias straight into plane copper. Whether that copper is
then coupled to a heatsink and whether forced air is present is NOT STATED - see Q6.
\item Duty cycle (continuous 200 W vs pulsed/keyed) is NOT STATED and materially changes the thermal
design - see Q8.
\end{itemize}
\subsubsection*{5. Size \& mounting}
NOT STATED in the brief - no outline, no mounting holes, no height limit given. See Q3 and Q4.

Constraints that are already fixed and will bound whatever outline is chosen:

\begin{itemize}
\item 4-layer stackup is MANDATORY (non-negotiable at 20 MHz). L2 must be an unbroken GND plane
directly under the power loop.
\item Board area is a design variable, not just a cost item: the fallback for the output inductor is a
PCB air-core spiral, which trades board area for BOM cost and sourcing risk. A tight HARD outline
cap could remove that fallback.
\item Copper pours for the tank and power loop must be wide (skin depth in Cu at 20 MHz is 14.6 um;
1 oz = 35 um is only \textasciitilde{}2.4 skin depths), so conductor area cannot be shrunk to DC-current sizing.
\end{itemize}
WARNING: any outline cap answered as HARD binds permanently at P5 board\_init and cannot be relaxed later without restarting layout. Answer HARD only if it is a real physical constraint.

\subsubsection*{6. Quantity \& budget}
NOT STATED in the brief - see Q1 and Q2.

The only budget guidance given is directional: "minimise spend by reducing SCOPE, never by reducing quality." Known BOM anchors: EPC2019 \textasciitilde{}\$2.17, LMG1020YFFR \textasciitilde{}\$0.36.

\subsubsection*{7. Assembly}
\begin{itemize}
\item \textbf{100\% JLCPCB PCBA. Zero off-catalog parts.} Owner ruled on this explicitly and it is the
reason the topology is Class E rather than the originally requested Class DE. Every part must be
LCSC/JLC stocked. This is a HARD constraint.
\item No hand-soldered parts, no wound magnetics, no THT-only parts. (This is what rejected the gate
drive transformer, the push-pull output transformer, and Mini-Circuits/Coilcraft RF transformers.)
\item 4-layer PCB.
\item Assembly sides: NOT STATED. ASSUMED single-sided (top) assembly, because the EPC2019 is
bottom-cooled into plane copper and a clear bottom side is useful for heatsinking - confirm
via Q6.
\item The two connectors (I1, I2) and the DC input (I3) must be JLC-assemblable parts, which
constrains the connector choices in Q5 and Q10.
\end{itemize}
\subsubsection*{8. Compliance / safety flags}
All three of the following APPLY to this board.

1. \textbf{\textgreater{}30 V present (40 V bus, 142.5 V peak switch node).} The DC bus is 40 V and the drain node swings to a designed 142.5 V peak, with real ringing on top of that - the 1.4x derate against the EPC2019's 200 V rating is the entire margin. Consequences: creepage/clearance on the drain net and tank must be sized for \textgreater{}=215 V peak (tank L/C see \textasciitilde{}215 V pk); the 1.4x derate is protected by POWER LOOP AREA, not trace width, so layout is a safety item here, not just performance. Exposed high-voltage nodes on an unenclosed bench board are a shock/arc hazard. 2. \textbf{High current (\textgreater{}3 A): \textasciitilde{}5.8 A DC input, \textasciitilde{}16 A peak switch current, 6.6 A rms tank current.} Conductor sizing must account for skin effect at 20 MHz (14.6 um skin depth), not DC ampacity. The output inductor dissipates \textasciitilde{}7 W and is the hottest component on the board - a burn hazard on an open bench board and a fire/derating risk if the chosen part is sized by current rating instead of dissipation. 3. \textbf{RF transmit: 200 W at 20 MHz.} This is a high-power HF RF source in a band adjacent to the 13.56 MHz and 27.12 MHz ISM allocations; 20 MHz itself is NOT an ISM allocation. EMC implications: a Class E stage is a hard-switched square-wave source, so harmonics at 40, 60, 80 MHz+ are strong and the only suppression is the output tank/match. Unenclosed, unshielded operation at 200 W will radiate. This board is for operation into a matched dummy load or a properly licensed/shielded setup only - it is not a certifiable transmitter and must not be connected to an antenna without an appropriate licence and additional filtering. 4. \textbf{Load-sensitivity hazard (consequence of the frozen topology).} Class E is narrowband and load-sensitive; it does not maintain ZVS across load pull. With protection logic explicitly out of scope, an open, shorted, or badly mismatched output at 200 W can drive Vds past the 1.4x margin and destroy the FET. This needs an explicit owner acknowledgement - see Q11.

Not applicable: no mains voltage anywhere on this board (external 40 V bench supply); no battery of any chemistry; no motors.

\subsubsection*{9. Open questions}
Answer these and the design can proceed. Defaults are offered where a sensible one exists - if a default is fine, just say "default".

1. \textbf{How many boards do you want built?} (Default: 5 - JLCPCB's typical minimum-cost quantity for a 4-layer assembled board.)

2. \textbf{What is your target cost per board, assembled?} A rough ceiling is enough (e.g. "under \$60 each"). This mainly decides whether the output inductor is bought (several paralleled SMD parts) or etched as a free PCB spiral that costs board area. (Default: no explicit ceiling - optimise for quality within a sensible bench-instrument cost.)

3. \textbf{Is there a maximum board size?} If yes, give width x height in mm, and say whether that limit is \textbf{HARD} (a real physical constraint, e.g. it must fit an existing fixture - this binds permanently and cannot be changed later) or \textbf{soft} (a preference). (Default: soft guidance of about 100 x 80 mm, no hard cap.)

4. \textbf{Do you want mounting holes?} If yes: how many, what screw size, and are their positions fixed by something existing? (Default: 4x M3 holes, one near each corner, positions chosen by the layout.)

5. \textbf{How should the 40 V DC bus connect to the board?} It must carry \textasciitilde{}6 A continuously. Options: (a) a screw terminal block, (b) a high-current 2-pin header, (c) solder lugs/pads for ring-terminal wires, (d) a banana-jack pair. Also: do you already have a 40 V supply that can deliver 6 A, and what does its output cable terminate in? (Default: (a) a 2-position screw terminal rated \textgreater{}=10 A.)

6. \textbf{How will the board be cooled?} It dissipates 35-50 W. Options: (a) bare board on the bench with no heatsink (this will get very hot and may not be survivable at continuous 200 W), (b) a heatsink bolted to the bottom of the board under the GaN FET, (c) a fan blowing across the top of the board, (d) both (b) and (c). Your answer decides whether the bottom copper must be left clear and flat, so it also fixes whether assembly is single- or double-sided. (Default: (d) - design for a bottom-side heatsink plus forced air, since that is the only answer that makes continuous 200 W comfortable.)

7. \textbf{What ambient (room) temperature should the board be designed to survive?} (Default: 25 C typical bench, designed with margin to 40 C.)

8. \textbf{Will this run at 200 W continuously, or in short bursts?} If bursts: roughly how long is a burst and how often (e.g. "10 seconds on, a minute off")? This changes the thermal design substantially - a pulsed duty cycle can make the hot output inductor a non-issue. (Default: continuous 200 W - the worst case, and the safe assumption.)

9. \textbf{Where should the 5 V driver supply come from?} Options: (a) an onboard step-down converter running off the 40 V bus, so the board needs only one power connection; (b) a separate 5 V input you feed from a bench supply or USB brick, which is one more cable but removes a part and removes 40 V switching noise near the gate drive. (Default: (a) onboard buck - one cable is worth more on a bench instrument.)

10. \textbf{Confirm SMA for the 200 W RF output.} The brief says SMA and SMA is electrically fine at 20 MHz, but SMA is a small connector for 200 W and repeated mating at that power is a known wear point; a BNC or N-type is more robust. (Default: keep SMA as briefed.)

11. \textbf{SAFETY - please confirm explicitly:} this board has NO protection of any kind (no VSWR detection, no overcurrent, no overvoltage, no thermal shutdown - all were placed out of scope). Class E is load-sensitive, so running it into an open, a short, or a badly mismatched load at 200 W can push the switch node past its voltage margin and destroy the FET, and the RF output is a burn/RF-exposure hazard at 200 W. Do you confirm this will only ever be operated into a proper 50 ohm dummy load or matched load, on a bench, by you? (No default - this must be answered.)

12. \textbf{What does your signal generator actually put out on the drive SMA?} Specifically: peak-to- peak amplitude, whether it can drive a 50 ohm load, and whether its duty cycle is adjustable at 20 MHz. This decides whether the LMG1020 can be driven directly or needs a bias network or a fast buffer in front of it. (Default: assume a 50 ohm-capable generator putting out \textasciitilde{}5 Vpp into 50 ohm with adjustable duty cycle, and design the input to tolerate 3.3-5 Vpp.)

---

\subsubsection*{10. ANSWERS - all 12 questions closed by the owner 2026-08-07 (P0 batch)}
\textbf{These are settled. Downstream phases treat them as requirements, not proposals.}

\begin{quote}\small\ttfamily\raggedright
\textbar{} \# \textbar{} Answer \textbar{}\\
\textbar{}---\textbar{}---\textbar{}\\
\textbar{} 1 \textbar{} Build quantity {*}{*}5{*}{*} (default taken) \textbar{}\\
\textbar{} 2 \textbar{} {*}{*}No hard cost ceiling{*}{*}; minimise spend by reducing scope, never quality (default taken) \textbar{}\\
\textbar{} 3 \textbar{} Outline {*}{*}\textasciitilde{}100 x 80 mm, SOFT - no hard cap{*}{*} (default taken). Deliberately NOT hard-capped so the PCB air-core spiral inductor fallback stays reachable \textbar{}\\
\textbar{} 4 \textbar{} {*}{*}4x M3{*}{*} mounting holes at corners (default taken) \textbar{}\\
\textbar{} 5 \textbar{} 40 V bus via {*}{*}2-position screw terminal, \textgreater{}=10 A{*}{*} (default taken) \textbar{}\\
\textbar{} 6 \textbar{} {*}{*}OWNER: continuous 200 W + bottom-side heatsink + forced air.{*}{*} Bottom copper stays clear and flat as a heatsink face -\textgreater{} {*}{*}single-sided (top) assembly{*}{*} \textbar{}\\
\textbar{} 7 \textbar{} Ambient {*}{*}25 C bench, margin to 40 C{*}{*} (default taken) \textbar{}\\
\textbar{} 8 \textbar{} {*}{*}OWNER: continuous 200 W.{*}{*} Worst case; thermal drives the layout \textbar{}\\
\textbar{} 9 \textbar{} {*}{*}5 V from an onboard buck{*}{*} off the 40 V bus - one cable (default taken) \textbar{}\\
\textbar{} 10 \textbar{} {*}{*}OWNER: keep SMA.{*}{*} 100 Vrms / 2 Arms at 20 MHz is within SMA HF limits; N-type would break the 100\% JLC PCBA rule \textbar{}\\
\textbar{} 11 \textbar{} {*}{*}OWNER SAFETY ACK: confirmed.{*}{*} No VSWR / OCP / OVP / thermal protection of any kind. Operated into a proper 50 ohm matched or dummy load ONLY. Owner accepts that an open, short or bad mismatch at 200 W will likely destroy the FET \textbar{}\\
\textbar{} 12 \textbar{} {*}{*}OWNER: \textasciitilde{}5 Vpp into 50 ohm, duty adjustable.{*}{*} -\textgreater{} 50 ohm termination direct into the LMG1020 input, {*}{*}NO buffer stage{*}{*}. Closes brief MUST-RESOLVE \#4 \textbar{}
\end{quote}
\paragraph*{Consequences that bind later phases}
\begin{itemize}
\item \textbf{Single-sided top assembly.} Bottom layer is reserved as a flat heatsink mounting face; no
bottom-side components. Thermal via array under the EPC2019 down to bottom copper.
\item \textbf{Thermal is a first-class layout driver}, not an afterthought: 35-50 W continuous, and the
output tank inductor (\textasciitilde{}7 W) runs hotter than the GaN.
\item \textbf{No protection parts} are to be added by any later phase. If a reviewer flags the missing
protection, the response is "waived, owner-acknowledged at P0" - not a new part.
\item \textbf{Outline stays soft} through P5 \texttt{board\_init}; do not pass a hard \texttt{--outline} cap that would
foreclose a PCB spiral inductor.
\end{itemize}
\section{Architecture}
\subsection*{Decisions of record (state.json)}
{\small
\begin{longtable}{lp{5.2cm}p{7.2cm}}
\toprule
\textbf{Phase} & \textbf{What} & \textbf{Why} \\
\midrule
\endhead
P0 & Thermal: continuous 200 W duty, bottom-side heatsink pad + thermal via array under EPC2019, forced air assumed; single-sided (top) assembly so bottom is a heatsink mounting face & Owner answer at P0 batch. 35-50 W loss with output inductor \textasciitilde{}7 W hotter than the GaN; worst-case design chosen \\
P0 & SAFETY ACK: no VSWR/OCP/OVP/thermal protection; matched or dummy 50 ohm load ONLY. Open/short output at 200 W will likely destroy the FET & Owner explicitly acknowledged. Class E loses ZVS under load pull; drain exceeds the 142.5 V design peak. Protection stays out of scope per power-stage-only brief \\
P0 & RF output connector = SMA (edge/PCB mount, 50 ohm, LCSC-stocked) & 200 W into 50 ohm at 20 MHz = 100 Vrms / 2 Arms, within SMA HF limits. N-type would break the 100\% JLC PCBA rule \\
P0 & Drive input: 50 ohm termination direct to LMG1020 input, NO buffer stage & Owner gen supplies \textasciitilde{}5 Vpp into 50 ohm with adjustable duty. Closes brief MUST-RESOLVE \#4; fewest parts and lowest added delay \\
P0 & Defaults taken (owner did not specify): qty 5; no hard cost ceiling; outline \textasciitilde{}100x80 mm SOFT (no hard cap); 4x M3 corner holes; 40 V bus via 2-pos screw terminal \textgreater{}=10 A; ambient 25 C bench w/ margin to 40 C; 5 V rail from onboard buck off the 40 V bus & Routine judgement calls. Outline deliberately left SOFT so the PCB air-core spiral inductor fallback stays available for the 6.6 A rms tank \\
P1 & CORRECTED EPC2019 specs from datasheet: Coss(er) 156 pF (not 110), package 2.77x0.95mm solder bars (not 1.35mm LGA), Rds(on) 22 typ/42 max (not 36), Qg 2.4/2.9 nC & Brief figures came from a web summary, not the datasheet. Two P1 agents independently verified. External C\_shunt drops to \textasciitilde{}161 pF \\
P1 & L\_s and L\_m to be etched as PCB air-core spirals, NOT catalog inductors & No LCSC part closes it: power inductors give Q 20-40 at 20 MHz (25-50 W); RF inductors rated 120-140 mA vs 6.6 A needed. Paralleling does NOT reduce loss (Q\_total = Q\_each). Spiral gives Q\textasciitilde{}130-150 at zero BOM cost \\
P1 & Two-zone floorplan mandatory: heatsink zone (Q1/L1/bulk/driver) vs magnetics zone (L\_s/L\_m/tank caps) with L2 plane keepout under spirals only & A metal heatsink or solid ground plane under a PCB spiral is a shorted turn. L2 stays unbroken under the power loop; cut out ONLY under spirals \\
\bottomrule
\end{longtable}
}
\subsection*{Sheet plan}
\subsection*{rf-de-20m - hierarchical sheet plan}
Three sheets under a root that contains nothing but the sheet symbols and the inter-sheet wiring. \textasciitilde{}58 parts. \textbf{Deliberately not decomposed further} - this is one function (a Class E stage) and the only organising principle that earns its keep is the floorplan, so the sheets map 1:1 onto the zones in \texttt{stackup.md} s3.

\begin{quote}\small\ttfamily\raggedright
\textbar{} Sheet \textbar{} Blocks (`blocks.md`) \textbar{} Zone \textbar{} Refdes block \textbar{} `\#PWR` base \textbar{}\\
\textbar{}---\textbar{}---\textbar{}---\textbar{}---\textbar{}---\textbar{}\\
\textbar{} {*}{*}`hk`{*}{*} \textbar{} B1 bus entry + bulk, B2 +5 V buck \textbar{} A \textbar{} {*}{*}100-199{*}{*} \textbar{} {*}{*}100{*}{*} \textbar{}\\
\textbar{} {*}{*}`stage`{*}{*} \textbar{} B3 drive input, B4 gate driver, B5 FET pair, B6 drain feed + bus decoupling \textbar{} A \textbar{} {*}{*}200-299{*}{*} \textbar{} {*}{*}200{*}{*} \textbar{}\\
\textbar{} {*}{*}`tank`{*}{*} \textbar{} B7 series tank, B8 L-match, B9 RF output \textbar{} B and C \textbar{} {*}{*}300-399{*}{*} \textbar{} {*}{*}300{*}{*} \textbar{}
\end{quote}
Refdes are unique across sheets by construction: each sheet owns a hundreds block for every prefix, and each sheet has its own \texttt{\#PWR} base.

---

\subsubsection*{1. Net-naming contract - BINDING ON P4}
KiCad names a net from a hierarchical label alone as \textbf{`/\textless{}sheet\textgreater{}/\textless{}LABEL\textgreater{}`}. A net that lives entirely inside one sheet therefore comes out with the sheet prefix, and \texttt{constraints.json} uses those names. {*}{*}A net whose name does not match is silently no-op'd by every P5-P8 consumer{*}{*} - this exact trap cost the lumina-par run a P4 amendment.

Power nets are \textbf{bare global symbols}: \texttt{+40V}, \texttt{+5V}, \texttt{GND}.

\textbf{Exactly one signal net crosses a sheet boundary: `/SW`.} For it to be named \texttt{/SW} and not \texttt{/stage/SW}, {*}{*}P4 must place a root-sheet local label spelled \texttt{SW} on the inter-sheet wire.{*}{*} Everything else is sheet-internal and is named with its sheet prefix on purpose.

\begin{quote}\small\ttfamily\raggedright
\textbar{} Canonical net \textbar{} Scope \textbar{} Carries \textbar{}\\
\textbar{}---\textbar{}---\textbar{}---\textbar{}\\
\textbar{} `+40V` \textbar{} global \textbar{} 5.94 A DC nominal, 7.0 A declared. Ring to 51 V \textbar{}\\
\textbar{} `+5V` \textbar{} global \textbar{} 143 mA avg, \textasciitilde{}1.6 A/gate peaks \textbar{}\\
\textbar{} `GND` \textbar{} global \textbar{} In1 + In2 + B.Cu plane stack, zones A and C \textbar{}\\
\textbar{} {*}{*}`/SW`{*}{*} \textbar{} {*}{*}`stage` -\textgreater{} `tank`, root label required{*}{*} \textbar{} drain node. {*}{*}9.0 A rms, 142.5 V pk{*}{*} \textbar{}\\
\textbar{} `/stage/DRIVE` \textbar{} `stage` \textbar{} J201 -\textgreater{} R201/R202 -\textgreater{} U201 IN+ \textbar{}\\
\textbar{} `/stage/GATE\_ON` \textbar{} `stage` \textbar{} U201 OUTH -\textgreater{} the two turn-on legs \textbar{}\\
\textbar{} `/stage/GATE\_OFF` \textbar{} `stage` \textbar{} U201 OUTL -\textgreater{} the two turn-off legs \textbar{}\\
\textbar{} `/stage/GATE\_Q1` \textbar{} `stage` \textbar{} gate of Q201 \textbar{}\\
\textbar{} `/stage/GATE\_Q2` \textbar{} `stage` \textbar{} gate of Q202 \textbar{}\\
\textbar{} {*}{*}`/tank/TANK\_A`{*}{*} \textbar{} `tank` \textbar{} L301 / C\_s junction. {*}{*}6.58 A rms, 170 V pk - the highest node on the board{*}{*} \textbar{}\\
\textbar{} `/tank/TANK\_B` \textbar{} `tank` \textbar{} C\_s / L302 junction. 6.58 A rms, 43 V pk \textbar{}\\
\textbar{} `/tank/RFOUT` \textbar{} `tank` \textbar{} L302 -\textgreater{} C\_m -\textgreater{} J301. 6.58 A rms at the C\_m node, 2.0 A into the load, 141 V pk \textbar{}
\end{quote}
\textbf{The OUTH/OUTL nets are named `\_ON`/`\_OFF`, NOT `\_H`/`\_L`, on purpose.} \texttt{rules\_gen.detect\_diff\_pairs} auto-pairs \texttt{high\_speed} nets whose names end in \texttt{\_H}/\texttt{\_L} (as well as \texttt{\_P}/\texttt{\_N}, \texttt{\_DP}/\texttt{\_DM}, \texttt{+}/\texttt{-}) and would silently treat the turn-on and turn-off legs as a 100 ohm differential pair, emitting a \texttt{diff\_pair\_gap} rule and an inner-layer \texttt{disallow track} rule on the two most inductance-critical nets on the board. Do not rename them back.

{*}\emph{Do not "tidy" `/stage/}\texttt{ or }/tank/\emph{` into bare names.}{*} Promoting a sheet-internal net to a root-crossing name means the root sees wire + label + one sheet pin, which raises ERC \texttt{label\_dangling}; buying the bare name would mean adding a part to the root sheet. The netlist is the authority.

---

\subsubsection*{2. \texttt{hk} (100-199, \texttt{\#PWR} 100) - zone A}
\begin{quote}\small\ttfamily\raggedright
\textbar{} Ref \textbar{} Part / class \textbar{} Note \textbar{}\\
\textbar{}---\textbar{}---\textbar{}---\textbar{}\\
\textbar{} {*}{*}J101{*}{*} \textbar{} 2-pos screw terminal, 5.08 mm, \textgreater{}= 24 A / 250 V (KF128-5.08-2P class) \textbar{} {*}{*}The only THT part on the board.{*}{*} Left edge, x \textless{} 5 mm, clear of the heatsink land (HS-3) \textbar{}\\
\textbar{} C101, C102 \textbar{} 2x 100 uF / {*}{*}63 V{*}{*} SMD polymer can \textbar{} bulk. SMD can, not leaded - no bottom-face solder \textbar{}\\
\textbar{} C103, C104 \textbar{} 2x 2.2 uF / 100 V {*}{*}X7R{*}{*} 0805 \textbar{} mid tier, 100 kHz-5 MHz. {*}{*}X7R is correct here and nowhere else on this board{*}{*} \textbar{}\\
\textbar{} {*}{*}U101{*}{*} \textbar{} {*}{*}LM5017-class 100 V synchronous COT buck{*}{*}, SOIC-8-EP \textbar{} Vin abs max \textgreater{}= 63 V; no comp network, no catch diode \textbar{}\\
\textbar{} L101 \textbar{} 15-47 uH shielded, \textgreater{}= 0.5 A \textbar{} placeholder value; {*}{*}P4 sizes against the chosen Fsw{*}{*} \textbar{}\\
\textbar{} C105 \textbar{} Vin local, 100 V \textbar{} \textbar{}\\
\textbar{} C106, C107 \textbar{} VCC decoupler, bootstrap \textbar{} per the buck datasheet \textbar{}\\
\textbar{} C108, C109 \textbar{} Vout bulk + HF on `+5V` \textbar{} \textbar{}\\
\textbar{} R101, R102 \textbar{} feedback divider \textbar{} sets 5.0 V +/-4\% \textbar{}\\
\textbar{} R103 \textbar{} RON / on-time resistor \textbar{} sets Fsw \textbar{}
\end{quote}
\textbf{Placement rule:} at the DC-input end, as far from the gate loop and the tank as zone A allows. Its 0.5-2 MHz switching noise is spectrally clear of 20 MHz but its switch node is still a dV/dt source next to a 1.4 V gate threshold.

---

\subsubsection*{3. \texttt{stage} (200-299, \texttt{\#PWR} 200) - zone A}
This sheet is the board. Everything on it lives in one \textasciitilde{}15 x 15 mm cluster except the drive connector.

\begin{quote}\small\ttfamily\raggedright
\textbar{} Ref \textbar{} Part / class \textbar{} Note \textbar{}\\
\textbar{}---\textbar{}---\textbar{}---\textbar{}\\
\textbar{} {*}{*}J201{*}{*} \textbar{} SMD edge-launch SMA, 50 ohm (BWSMA-KE-P001 class) \textbar{} {*}{*}SMD, not the THT bulkhead jack{*}{*} - a THT jack solders through the bottom face and breaks the heatsink land \textbar{}\\
\textbar{} R201, R202 \textbar{} {*}{*}2x 100 R 0805 in parallel{*}{*} = 50 ohm \textbar{} \textasciitilde{}0.125 W in the termination; a single 0603 is over its rating. {*}{*}DC-coupled - see `blocks.md` B3{*}{*} \textbar{}\\
\textbar{} {*}{*}U201{*}{*} \textbar{} {*}{*}LMG1020YFFR{*}{*}, 6-ball WCSP 0.8 x 1.2 mm \textbar{} {*}{*}Sits ON the mirror axis, equidistant from both gate bars.{*}{*} LCSC brands it "Tokmas" - authenticity check on receipt \textbar{}\\
\textbar{} C201 \textbar{} 100 nF 0402 X7R \textbar{} VDD bypass, {*}{*}\textless{} 0.5 mm from the VDD ball, via-in-pad return{*}{*} \textbar{}\\
\textbar{} C202 \textbar{} 10 nF 0201 C0G \textbar{} second VDD bypass, straddling the ball \textbar{}\\
\textbar{} R203, R204 \textbar{} 2x 4.0 R 0603 in parallel \textbar{} OUTH -\textgreater{} `/stage/GATE\_Q1` \textbar{}\\
\textbar{} R205, R206 \textbar{} 2x 4.0 R 0603 in parallel \textbar{} OUTH -\textgreater{} `/stage/GATE\_Q2` \textbar{}\\
\textbar{} R207, R208 \textbar{} 2x 4.0 R 0603 in parallel \textbar{} OUTL -\textgreater{} `/stage/GATE\_Q1` \textbar{}\\
\textbar{} R209, R210 \textbar{} 2x 4.0 R 0603 in parallel \textbar{} OUTL -\textgreater{} `/stage/GATE\_Q2` \textbar{}\\
\textbar{} {*}{*}Q201, Q202{*}{*} \textbar{} {*}{*}2x EPC2019{*}{*} 200 V eGaN \textbar{} {*}{*}MIRRORED PAIR about the U201 axis.{*}{*} Both populated by default \textbar{}\\
\textbar{} C203-C206 \textbar{} 4x 56 pF / 1 kV C0G 1206 \textbar{} {*}{*}C\_shunt TRIM - DNP by default.{*}{*} In the power loop, not on a stub \textbar{}\\
\textbar{} {*}{*}L201{*}{*} \textbar{} RF choke, {*}{*}\textgreater{}= 0.82 uH, SRF \textgreater{}= 80 MHz, DCR \textless{}= 25 mohm, I\_sat \textgreater{}= 12 A, \textgreater{}= 8 A rms{*}{*} \textbar{} Hardest part on the BOM. {*}{*}Pre-authorised escape: 2x 0.47 uH in series.{*}{*} Not authorised: a smaller choke - see `blocks.md` B6 \textbar{}\\
\textbar{} C207-C210 \textbar{} 4x 10 nF / 100 V C0G 0603 \textbar{} bus HF bank, {*}{*}within 3 mm of L201's bus-side pad{*}{*} \textbar{}\\
\textbar{} C211, C212 \textbar{} 2x 1 nF / 100 V C0G 0603 \textbar{} terminates the bank's self-resonance \textbar{}
\end{quote}
\textbf{Four resistors per FET pair per polarity is not over-specification.} Individual gate resistors are what damp the differential mode between the two gate loops; a shared resistor leaves the two gate nodes coupled through the driver output and free to oscillate. Two 0603s in parallel per leg rather than one 0805 keeps each part at 0.047 W (a single part is at or over its derated rating at 100 C) and roughly halves the leg's parasitic inductance. See \texttt{blocks.md} s3.

{*}{*}The gate-loop budget is the acceptance criterion for this whole cluster: \textless{}= 0.8 nH per FET, matched to +/-0.1 nH.{*}{*} At R = 3.1 ohm that mismatch is 32 ps of skew against a \textasciitilde{}1 ns edge. {*}{*}Match the loops geometrically, not the trace lengths electrically{*}{*} - propagation on FR4 is \textasciitilde{}6.7 ps/mm, so length matching is not the mechanism.

\textbf{Common source star.} Both source bars return to \textbf{one} via cluster on the mirror axis, and the LMG1020's GND connects to the source plane \textbf{only at that star point} (TI LMG1020DS s10.1.1 - single-point ground). Separate source returns would make common-source inductance differential and imbalance the pair.

---

\subsubsection*{4. \texttt{tank} (300-399, \texttt{\#PWR} 300) - zones B and C}
\begin{quote}\small\ttfamily\raggedright
\textbar{} Ref \textbar{} Part / class \textbar{} Zone \textbar{} Note \textbar{}\\
\textbar{}---\textbar{}---\textbar{}---\textbar{}---\textbar{}\\
\textbar{} {*}{*}L301{*}{*} \textbar{} {*}{*}L\_s 184 nH - etched PCB air-core spiral{*}{*} \textbar{} B \textbar{} 2 turns, OD 30-34 mm, 1.5 mm trace, \textasciitilde{}900-1150 mm\textasciicircum{}2. {*}{*}No LCSC code. `PCB feature - do not place`.{*}{*} SPIRAL-1..6 \textbar{}\\
\textbar{} C301-C308 \textbar{} 8x 56 pF / {*}{*}1 kV{*}{*} C0G 1206 = 448 pF \textbar{} B \textbar{} C\_s. 0.82 A rms and \textasciitilde{}145 mW each. {*}{*}No X7R{*}{*} \textbar{}\\
\textbar{} {*}{*}L302{*}{*} \textbar{} {*}{*}L\_m 115 nH - etched PCB air-core spiral{*}{*} \textbar{} B \textbar{} {*}{*}\textgreater{}= 900 mm\textasciicircum{}2 despite being the smaller inductance{*}{*} - the constraint is dissipation area, not L (SPIRAL-3) \textbar{}\\
\textbar{} C309-C317 \textbar{} 9x 56 pF / 1 kV C0G 1206 = \textasciitilde{}504 pF \textbar{} C \textbar{} C\_m. Carries {*}{*}6.27 A rms{*}{*}, not the 2 A load current - 0.70 A each \textbar{}\\
\textbar{} {*}{*}J301{*}{*} \textbar{} SMD edge-launch SMA, 50 ohm - {*}{*}same part as J201{*}{*} \textbar{} C \textbar{} right edge. 100 Vrms / 2.0 A rms \textbar{}
\end{quote}
\textbf{Why C\_m sits in zone C and C\_s does not.} C\_s is a \emph{series} element and floats - it needs no ground reference and belongs beside the coils. C\_m returns to GND and carries 6.27 A rms, so it must sit where the ground plane exists.

\textbf{L301 and L302 are schematic symbols with custom footprints, not annotations.} The footprint carries the spiral copper on F.Cu (and B.Cu under SPIRAL-2), a courtyard equal to the thermal footprint, and its two terminal pads. Copper that no tool knows about gets routed over, poured under and placed on top of.

---

\subsubsection*{5. Placement plan for P6}
\paragraph*{5.1 Placement groups (declared in \texttt{constraints.json.placement.groups})}
\begin{quote}\small\ttfamily\raggedright
\textbar{} Group \textbar{} Anchor \textbar{} Members \textbar{} Purpose \textbar{}\\
\textbar{}---\textbar{}---\textbar{}---\textbar{}---\textbar{}\\
\textbar{} `switch` \textbar{} {*}{*}Q201{*}{*} \textbar{} Q202, U201, R203-R210, C201-C212, L201 \textbar{} {*}{*}The whole thing.{*}{*} Power loop + gate loop + choke + bus HF in one \textasciitilde{}15 x 15 mm cluster \textbar{}\\
\textbar{} `drive\_in` \textbar{} J201 \textbar{} R201, R202 \textbar{} 50 ohm termination at the connector, not at the driver \textbar{}\\
\textbar{} `tank\_ls` \textbar{} L301 \textbar{} C301-C308 \textbar{} C\_s bank adjacent to the L\_s spiral \textbar{}\\
\textbar{} `tank\_lm` \textbar{} L302 \textbar{} C309-C317, J301 \textbar{} L\_m -\textgreater{} C\_m -\textgreater{} SMA, in that order, shortest possible \textbar{}\\
\textbar{} `hk` \textbar{} U101 \textbar{} L101, C105-C109, R101-R103 \textbar{} buck and its loop \textbar{}\\
\textbar{} `bus\_in` \textbar{} J101 \textbar{} C101-C104 \textbar{} bulk at the terminal \textbar{}
\end{quote}
\paragraph*{5.2 Edges}
\begin{quote}\small\ttfamily\raggedright
\textbar{} Ref \textbar{} Edge \textbar{} pos \textbar{} Why \textbar{}\\
\textbar{}---\textbar{}---\textbar{}---\textbar{}---\textbar{}\\
\textbar{} J101 \textbar{} {*}{*}left{*}{*} \textbar{} 0.5 \textbar{} DC entry. {*}{*}x \textless{} 5 mm so its THT pins clear the heatsink land{*}{*} (HS-3) \textbar{}\\
\textbar{} J201 \textbar{} {*}{*}top{*}{*} \textbar{} 0.2 \textbar{} drive input, in zone A near the driver \textbar{}\\
\textbar{} J301 \textbar{} {*}{*}right{*}{*} \textbar{} 0.5 \textbar{} RF output, in zone C. Must be adjacent to the C\_m node \textbar{}
\end{quote}
\paragraph*{5.3 Parts that must be placed by explicit \texttt{place\_edit} and LOCKED before the annealer runs}
The annealer optimises a cost function; none of the constraints below are expressible in it.

\begin{quote}\small\ttfamily\raggedright
\textbar{} Ref \textbar{} Why locking is mandatory \textbar{}\\
\textbar{}---\textbar{}---\textbar{}\\
\textbar{} {*}{*}Q201, Q202{*}{*} \textbar{} mirror symmetry about the U201 axis is the whole of `blocks.md` s4.1(c). An annealer will not produce a mirror pair \textbar{}\\
\textbar{} {*}{*}U201{*}{*} \textbar{} must sit {*}{*}on{*}{*} the axis, equidistant from both gate bars \textbar{}\\
\textbar{} {*}{*}R203-R210{*}{*} \textbar{} four matched legs; two arms must be geometric mirrors \textbar{}\\
\textbar{} {*}{*}L301, L302{*}{*} \textbar{} zone B, centre-to-centre \textgreater{}= 38 mm, \textgreater{}= 14 mm clear of the heatsink land (SPIRAL-5, SPIRAL-6) \textbar{}\\
\textbar{} {*}{*}J301{*}{*} \textbar{} zone C, adjacent to the C\_m node, so `/tank/RFOUT` stays \textless{}= 15 mm \textbar{}\\
\textbar{} {*}{*}J101{*}{*} \textbar{} x \textless{} 5 mm (HS-3) \textbar{}
\end{quote}
{*}{*}After L301 and L302 are placed and locked, hand-add KiCad rule areas (keepouts) over both courtyards on all four layers.{*}{*} No pipeline check enforces "no other copper here" and \texttt{planes\_gen} has no void support. Verify geometrically at P8.

\paragraph*{5.4 Zone assignment for P6 (board-local, translate by \texttt{outline\_bbox} first)}
\begin{quote}\small\ttfamily\raggedright
\textbar{} Zone \textbar{} x \textbar{} Sheets \textbar{} Contents \textbar{}\\
\textbar{}---\textbar{}---\textbar{}---\textbar{}---\textbar{}\\
\textbar{} {*}{*}A{*}{*} \textbar{} 0 - 48 \textbar{} `hk`, `stage` \textbar{} everything except the tank \textbar{}\\
\textbar{} {*}{*}B{*}{*} \textbar{} 48 - 88 \textbar{} `tank` \textbar{} L301, L302, C301-C308. {*}{*}Nothing else may be placed here{*}{*} \textbar{}\\
\textbar{} {*}{*}C{*}{*} \textbar{} 88 - 100 \textbar{} `tank` \textbar{} C309-C317, J301 \textbar{}
\end{quote}
Heatsink land on B.Cu: \textbf{{[}5, 10, 36, 70{]}}, declared as a back-side keepout.

---

\subsubsection*{6. Sheet-level P4 notes}
1. \textbf{The C\_shunt trim bank (C203-C206) must be in the netlist}, marked DNP. A DNP part still has pads, so its clearance and area are accounted for at P6/P7 and the one-FET fallback stays a populate change rather than a respin. Mark them \texttt{DNP} in the BOM variant field so \texttt{bom\_cpl} excludes them at P9. 2. \textbf{`/SW` needs a root-sheet local label} or it becomes \texttt{/stage/SW} and six \texttt{constraints.json} entries silently stop matching. 3. \textbf{Both spirals need schematic symbols.} Do not let them become a layout-only feature; their nets must be real. 4. \textbf{The drive input is DC-coupled.} Do not add a series blocking capacitor "for safety" - it destroys duty-cycle control (\texttt{blocks.md} B3). Record it on the schematic. 5. \textbf{No protection parts.} No TVS, no fuse, no clamp, no OCP, no OVP, no thermal shutdown. Owner-acknowledged at P0 Q11. If ERC or a reviewer wants one, the response is \emph{"waived, owner-acknowledged at P0"}. 6. \textbf{L201's value is a floor, not a target.} \textgreater{}= 0.82 uH. If P3 cannot source it, the answer is 2x 0.47 uH in series - \textbf{not} a smaller single part. 7. \textbf{`IN-` on U201 ties to GND.} No bias divider, no buffer. 8. \textbf{Silkscreen hazard markings}: zone B (two \textasciitilde{}130 C copper structures, one on each face if SPIRAL-2 is adopted), \texttt{/tank/TANK\_A} (170 V pk), \texttt{/SW} (142.5 V pk), and J301 (200 W RF - burn and RF-exposure hazard). This is an unenclosed bench board with no interlocks of any kind.

Full architecture narrative (not inlined; see the artifact index): \texttt{architecture/blocks.md}, \texttt{architecture/decisions.md}, \texttt{architecture/power\_tree.md}, \texttt{architecture/stackup.md}.
\section{Schematic}
\emph{Pending --- produced at P4.}
\section{Layout}
\emph{Pending --- produced at P6.}
\section{Verification}
\emph{Pending --- produced at P8.}
\section{DFM and Fabrication}
\emph{Pending --- produced at P9.}
\section{Run Record (Appendix)}
\subsection*{Phase digests}
\paragraph*{P0}
\emph{(no digest recorded)}
\paragraph*{P1}
\emph{(no digest recorded)}
\paragraph*{P2}
\begin{quote}\small\ttfamily\raggedright
\# P2 Architecture digest - rf-de-20m (20 MHz Class E GaN stage, 200 W)\\
~\\
Artifacts: `architecture/blocks.md`, `power\_tree.md`, `stackup.md`, `sheets.md`,\\
`constraints.json`, `decisions.md`.\\
~\\
- {*}{*}Chain:{*}{*} SMA drive in -\textgreater{} 50 ohm DC-coupled termination -\textgreater{} LMG1020 -\textgreater{} {*}{*}two\\
  paralleled EPC2019 in a mirrored pair{*}{*} -\textgreater{} Coss-only shunt (312 pF) -\textgreater{} etched PCB\\
  spiral L\_s 184 nH -\textgreater{} 8x 56 pF 1 kV C0G -\textgreater{} etched spiral L\_m 115 nH -\textgreater{} 9x 56 pF\\
  C\_m -\textgreater{} SMA out. Housekeeping: LM5017-class 100 V COT buck off the 40 V bus.\\
- {*}{*}Stackup:{*}{*} `JLC04161H-1080B`, 4L, 1.6 mm, 1 oz outer / 0.5 oz inner. Chosen for\\
  {*}{*}one number - the 0.2444 mm L1-L2 dielectric{*}{*}, which is what makes EPC's stacked\\
  self-cancelling power loop possible. Not impedance, not density. {*}{*}1 oz, not 2 oz{*}{*}:\\
  the 2 oz lamination doubles the L1-L2 gap and therefore the power-loop inductance.\\
- {*}{*}Fab class:{*}{*} 100 x 80 mm, {*}{*}TG155+ high-Tg{*}{*} (spiral copper 100-140 C), {*}{*}ENIG\\
  not HASL{*}{*} (EPC AN009; solder-bar LGA + 0.4 mm WCSP), {*}{*}epoxy-filled-and-capped\\
  via-in-pad{*}{*} (LMG1020 VDD/GND, without which the VDD loop exceeds 0.3 nH). All\\
  three are order-time options, none expressible in `stackups.yaml`.\\
- {*}{*}Cost:{*}{*} \textasciitilde{}\$28-35 per assembled board at qty 5 (\textasciitilde{}\$140-175 for 5), plus \$10-20 of\\
  off-board heatsink and fan. No cost ceiling was set (P0 Q2).\\
- {*}{*}`constraints.json` lint-clean{*}{*}: 0 errors, 0 warnings (4 high\_speed, 6 power,\\
  explicit empty diff\_pairs, 5 voltages, {*}{*}3 voltage\_pairs{*}{*}, 7 thermal, 6 planes,\\
  placement edges/groups/keepouts/separation/fixed).\\
~\\
\#\# The four rulings\\
~\\
1. {*}{*}THE THERMAL BLOCKER -\textgreater{} build for TWO EPC2019.{*}{*} RthJB 7.5 C/W is inside the\\
   package: one FET is Tj {*}{*}138 C nominal{*}{*} and {*}{*}\textasciitilde{}203 C on a max-Rds part out of\\
   JLC's reel{*}{*} - not a viable build. Two dies: Tj {*}{*}103 C nominal, \textasciitilde{}141 C worst{*}{*}\\
   (under the 150 C abs max). {*}{*}k = 2 is simultaneously the ZVS ceiling and the\\
   thermal optimum{*}{*}: `Coss(er) \textless{}= 317 pF` caps the area factor at 2.03, and\\
   2 x 156 = 312 pF, so {*}no single larger LCSC GaN can beat the pair - only tie it{*}.\\
   Cost: +\$2.17, {*}{*}-0.9 pt efficiency{*}{*}, and a mirror-symmetry layout discipline.\\
   Bought: -35 C, a survivable worst case, no external shunt cap, and the heatsink\\
   relaxed {*}{*}0.3 -\textgreater{} 1.4 C/W{*}{*} - which is what makes the two-zone floorplan\\
   geometrically possible at all.\\
2. {*}{*}C\_shunt -\textgreater{} the datasheet wins.{*}{*} `output-network.json`'s 4 x 56 pF = 224 pF is\\
   {*}{*}withdrawn{*}{*} (derived from the brief's refuted 110 pF Coss; 39\% high, would lose\\
   ZVS). With the pair there is {*}{*}no external shunt cap{*}{*}; C203-C206 become four DNP\\
   trim sites in the power loop.\\
3. {*}{*}L\_s and L\_m are etched spirals, as first-class pipeline objects{*}{*} - BOM line\\
   with no LCSC code, schematic symbol, custom footprint, explicit place + lock,\\
   hand-added four-layer keepout. Six acceptance criteria SPIRAL-1..6.\\
4. {*}{*}Two-zone floorplan -\textgreater{} the conflict was an artefact.{*}{*} The power loop {*}ends{*} at\\
   the drain node and the tank {*}begins{*} there, so "unbroken L2 under the power loop"\\
   and "no plane under the spirals" apply to disjoint regions. Three bands:\\
   A {[}0,48{]} full plane stack + heatsink, B {[}48,88{]} {*}{*}no inner or bottom copper, no\\
   heatsink{*}{*}, C {[}88,100{]} planes restored for the RF output.\\
~\\
\#\# Findings no P1 fragment had\\
~\\
- {*}{*}`/tank/TANK\_A` is the highest node on the board at 170 V pk - 27 V above the\\
  drain.{*}{*} Series-resonant magnification at Q\_L = 5. The brief's "\textasciitilde{}215 V pk on the\\
  tank" is the voltage {*}across L\_s{*}, not a node voltage; the genuine element-across\\
  differentials (215 / 165 / 135 V pk) go into {*}{*}`voltage\_pairs`{*}{*}, where node\\
  arithmetic would give 57 V for the L\_s pair - 3.8x too low.\\
- {*}{*}The gate-loop spec is 0.8 nH per FET, not 0.3 nH.{*}{*} `power.md`'s 0.3 nH assumes\\
  a zero-resistance loop; with TI's mandatory \textgreater{}= 2 ohm resistor the peak gate current\\
  is 1.6 A, not 7 A, and EPC WP008 Eq.1 gives `L\_G \textless{}= 0.84 nH`. Matching is\\
  {*}{*}+/-0.1 nH = 32 ps of L/R skew{*}{*} - a {*}loop{*} problem, not a length-matching\\
  problem (FR4 is 6.7 ps/mm, so 15 mm of length error is only 100 ps).\\
- {*}{*}AC-coupling the drive input would break duty-cycle control.{*}{*} DC restore pins the\\
  average at the bias point; at D = 0.4 both logic levels sit above threshold and the\\
  FETs never turn off. -\textgreater{} DC-coupled, generator must be unipolar 0/+5 V ({*}{*}OPEN-1{*}{*}).\\
- {*}{*}SPIRAL-2 (verify-later upside):{*}{*} paralleling each spiral's winding on F.Cu AND\\
  B.Cu halves its resistance for the same L - magnetics loss {*}{*}16.3 -\textgreater{} \textasciitilde{}10 W, +2 pts\\
  of board efficiency for zero BOM cost.{*}{*} Not banked; every number in\\
  `power\_tree.md` is quoted for the conservative single-layer case.\\
- {*}{*}L\_m needs a spiral nearly as large as L\_s{*}{*} despite being 62\% of the inductance,\\
  because the binding constraint is dissipation area, not L.\\
~\\
\#\# Pipeline traps - two carried, one retired\\
~\\
1. {*}{*}`rules\_gen` only solves impedance for DIFFERENTIAL pairs{*}{*} (read\\
   `detect\_diff\_pairs` / `diff\_pair\_rules`). A lone single-ended `high\_speed` entry\\
   with `impedance\_ohm: 50` emits {*}{*}no width rule at all{*}{*} - it looks declared and\\
   does nothing. `/tank/RFOUT` is therefore declared under `power`, which is also the\\
   right physics: at 20 MHz a 15 mm run is lambda/550 and perturbs the load \textless{}2\%,\\
   while the 0.4332 mm `se\_50` geometry would run 50-75 C hot at 2 A rms. {*}{*}This\\
   contradicts a frozen requirement and is raised as OPEN-2, not absorbed.{*}{*}\\
2. {*}{*}`detect\_diff\_pairs` auto-pairs `high\_speed` nets ending `\_H`/`\_L`.{*}{*} The\\
   OUTH/OUTL legs are named `/stage/GATE\_ON` / `\_OFF` on purpose - `\_H`/`\_L` would\\
   have put a `diff\_pair\_gap` rule and an inner-layer `disallow track` rule on the\\
   two most inductance-critical nets on the board.\\
3. {*}{*}`planes\_gen` still has no void or keepout support{*}{*}, so zone B is unpoured by\\
   construction (six region entries, three layers, two regions each), and the spiral\\
   courtyards need hand-added KiCad rule areas at P6 plus geometric verification at\\
   P8. `board\_init` still does not origin the outline at (0,0) - {*}{*}translate every\\
   rect by `outline\_bbox` at P5.{*}{*}\\
4. {*}{*}RETIRED:{*}{*} the lumina-run trap "`rules\_gen` puts every power net in ONE Power\\
   class at the widest width" is {*}{*}fixed in the current script{*}{*} - it now buckets one\\
   class per required width. Declaring `/tank/RFOUT` at 6.6 A no longer widens `+5V`.\\
~\\
\#\# Riskiest decision\\
~\\
{*}{*}Paralleling the FETs trades a well-characterised loss for a poorly-characterised\\
one.{*}{*} With two dies, {*}{*}Coss hysteresis becomes the dominant term (5.7 W of\\
10.9 W){*}{*} and it rests on an assumed 10\% of stored energy that {*}{*}EPC does not\\
publish{*}{*}; at 15\% the pair's loss is \textasciitilde{}12.7 W. The mitigation is already in the\\
architecture and costs nothing: the trim-site footprints and the fact that a\\
two-FET layout builds fine with one FET populated. {*}{*}First bring-up step: measure DC\\
input power at 200 W out and thermal-image both dies.{*}{*}\\
~\\
\#\# Sim candidates nominated for the P8 `sim` gate\\
~\\
SIM-1 `classe\_zvs\_nominal` (Vds at turn-on \textless{}= 5 V, dVds/dt within +/-0.2 V/ns,\\
Pout 190-210 W, Vds peak \textless{}= 155 V) - {*}{*}the bench that decides whether this board\\
works.{*}{*} SIM-2 `cshunt\_sweep` (280-420 pF; retune \textless{}= +/-15\% must restore SIM-1) -\\
{*}{*}must pass before P4 fixes any tank value.{*}{*} SIM-3 `gate\_symmetry` (at +/-0.15 nH\\
mismatch: Id imbalance \textless{}= 15\%, Vgs inside -4.0/+5.75 V). SIM-4 `tank\_match`\\
(Zin 4.15-5.08 ohm at 20 MHz, H2 \textless{}= -20 dBc, H3 \textless{}= -33 dBc, {*}{*}including the\\
spiral-to-spiral mutual k{*}{*}). SIM-5 `loss\_split` (FET pair \textless{}= 12.0 W typ /\\
\textless{}= 17.0 W max-Rds - the numbers the 1.4 C/W heatsink was solved against).\\
SIM-6 `gate\_loop` (advisory).
\end{quote}
\subsection*{Gate attempt history}
\emph{(no entries)}
\subsection*{Issues}
\emph{(no entries)}
\subsection*{Human checkpoints}
{\small
\begin{longtable}{lllp{2.6cm}p{6.4cm}}
\toprule
\textbf{Id} & \textbf{Phase} & \textbf{Status} & \textbf{When} & \textbf{Note} \\
\midrule
\endhead
1 & P2 & not reached &  &  \\
2 & P4 & not reached &  &  \\
3 & P6 & not reached &  &  \\
4 & P8 & not reached &  &  \\
5 & P10 & not reached &  &  \\
\bottomrule
\end{longtable}
}
\section{Artifact Index}
{\small
\begin{longtable}{p{9cm}rp{4cm}}
\toprule
\textbf{File} & \textbf{Size} & \textbf{Note} \\
\midrule
\endhead
\texttt{state.json} & 5,991 B &  \\
\texttt{brief/brief.md} & 4,153 B &  \\
\texttt{requirements.md} & 16,663 B &  \\
\texttt{architecture/blocks.md} & 33,156 B &  \\
\texttt{architecture/constraints.json} & 19,958 B &  \\
\texttt{architecture/decisions.md} & 26,772 B &  \\
\texttt{architecture/power\_tree.md} & 15,268 B &  \\
\texttt{architecture/sheets.md} & 12,601 B &  \\
\texttt{architecture/stackup.md} & 17,322 B &  \\
\bottomrule
\end{longtable}
}
\end{document}
